\documentclass[11pt]{article}

\usepackage[margin=1in]{geometry}
\usepackage[T1]{fontenc}
\usepackage[utf8]{inputenc}
\usepackage{lmodern}
\usepackage{graphicx}
\usepackage{hyperref}

\title{
\texttt{libframe} : Low Level Image Processing
}
\author{}
\date{}

\begin{document}

\maketitle

This writeup is intended to match the current goals of \texttt{libframe} as
described in the project README. The project is not only about calibrating
rolling-shutter webcams and an Arducam, but also about low-level camera
control, stereo rectification, depth-oriented processing, and experiments on
camera behavior through exposed controls such as auto exposure, white
balance, and gain. In that sense, the document focuses both on geometric
calibration and on low-level camera details that are relevant to 3A-style
analysis and capture stability.

\section{Rolling Shutter vs.\ Global Shutter}

A global shutter camera exposes all pixels at the same instant. Each frame
therefore corresponds to a single capture time, which matches the standard
pinhole-camera assumption used by most calibration methods. If the camera or
scene moves a little during exposure, the entire frame still represents one
geometric snapshot.

A rolling shutter camera does not expose all rows at once. Instead, the
sensor is read out line by line, so the top and bottom of the image
correspond to slightly different capture times. When the camera or the
calibration target moves during this readout, the image is no longer a
perfect perspective projection of a rigid scene at one instant. Straight
lines may appear bent or slanted, and a planar calibration board can look
geometrically warped.

This matters because standard calibration algorithms assume that all observed
corners in one image came from one fixed camera pose. That assumption is
approximately true for rolling shutter cameras only when motion during
capture is negligible. In practice, rolling shutter causes trouble when:

\begin{itemize}
\item the board is moved by hand while capturing,
\item the camera is handheld or vibrating,
\item the exposure time is long because the scene is dim,
\item stereo cameras are not synchronized closely enough, or
\item calibration images are taken during panning or rapid viewpoint changes.
\end{itemize}

For this project, the correct approach is to first calibrate the rolling
shutter webcams using the ordinary camera model, but to collect the data in
a way that minimizes rolling-shutter distortion. This gives a useful
intrinsic calibration for static scenes and is usually the right first step
before considering a more advanced rolling-shutter-aware model.

\section{What Camera Calibration Is}

Camera calibration estimates the parameters that map 3D scene points to 2D
image measurements.

The most important outputs are:

\begin{itemize}
\item \textbf{Intrinsic parameters:} focal lengths, principal point, and
      lens distortion coefficients.
\item \textbf{Extrinsic parameters:} the pose of the camera relative to the
      calibration target for each image.
\item \textbf{Stereo parameters:} for a stereo pair, the rigid transform
      between the left and right cameras, along with rectification maps.
\end{itemize}

In our case, calibration is needed so that later steps such as rectification,
disparity estimation, depth recovery, and feature triangulation use correct
geometry rather than raw distorted image coordinates.

\section{Procedure Used in This Project}

The repository uses a ChArUco board, which combines ArUco markers with a
chessboard layout. This gives two advantages: marker IDs help with robust
detection, and the interpolated chessboard corners provide more accurate
calibration points.

Our current procedure is:

\begin{enumerate}
\item Print or prepare one fixed ChArUco board with known parameters:
      number of squares in $x$ and $y$, square length, marker length, and
      dictionary.
\item Capture a set of images with \texttt{capture\_charuco}.
\item Detect ArUco markers in each image.
\item Interpolate ChArUco corners from the marker detections.
\item Use the known board geometry and the detected 2D corners to estimate
      the camera matrix and distortion coefficients.
\item Reject large-error outlier images and repeat calibration for a small
      number of passes.
\item Save the final intrinsic calibration to a JSON file.
\end{enumerate}

The calibration succeeds only if the board description used during
calibration matches the board used during capture. In particular,
\texttt{squares-x} and \texttt{squares-y} refer to the number of squares on
the board, not the number of inner corners. An incorrect board specification
can produce very large reprojection errors even if corner detection looks
successful.

\section{Practical Calibration Procedure for Rolling Shutter Webcams}

Because our webcams are rolling shutter devices, the calibration images must
be collected carefully:

\begin{itemize}
\item Mount the camera rigidly if possible.
\item Keep the ChArUco board still when each frame is captured.
\item Use bright lighting to keep exposure time short.
\item Avoid hand motion, panning, or fast viewpoint changes while taking
      images.
\item Capture a wide range of board poses: fronto-parallel views, tilted
      views, and views near the image edges and corners.
\item Remove blurry or obviously skewed images before calibration.
\end{itemize}

The goal is to make each image behave as closely as possible to a
global-shutter snapshot. Under these conditions, standard calibration is
usually good enough for static-scene stereo work.

\section{Camera Settings for Calibration}

Besides board motion and camera motion, camera control settings also affect
how repeatable calibration images are. A calibration sequence works best when
the appearance of the board stays consistent from frame to frame. Large
automatic changes in exposure or white balance can make corner localization
less stable and make the dataset harder to reason about.

For the webcam controls we inspected, the important settings are auto
exposure, exposure time, gain, white balance, anti-flicker mode, and
backlight compensation.

The current values are usable for rough capture, but they are not the most
stable choice for calibration:

\begin{itemize}
\item \texttt{white\_balance\_automatic=1} means auto white balance is still
      changing color response across frames.
\item \texttt{auto\_exposure=3} means the camera is still choosing exposure
      automatically.
\item \texttt{exposure\_time\_absolute} is inactive in that mode, so exposure
      cannot be fixed directly.
\item \texttt{gain=0} is good if the scene is bright enough, because low gain
      reduces noise.
\item \texttt{power\_line\_frequency=50 Hz} is only appropriate under 50 Hz
      lighting. For a webcam used in the United States, 60 Hz is usually the
      better anti-flicker choice.
\item \texttt{backlight\_compensation=80} may make exposure less predictable,
      so it is better set to zero for calibration if possible.
\end{itemize}

The recommended approach is to lock the camera as much as the device allows:

\begin{itemize}
\item disable auto white balance and set a fixed white balance temperature,
\item disable auto exposure and choose a fixed exposure time,
\item keep gain low,
\item set the power-line frequency to match the room lighting, and
\item disable backlight compensation.
\end{itemize}

In practice, this means the calibration operator should first make the scene
bright enough for a short exposure, then switch to manual settings and keep
those settings unchanged while collecting all calibration images. If manual
exposure produces dim images, the correct fix is usually more light rather
than more gain. This is especially helpful for rolling shutter webcams,
because brighter scenes allow shorter exposures and therefore reduce motion
artifacts during readout.

\section{Example of a Weak Calibration Frame}

Figures~\ref{fig:auto-frame} and \ref{fig:bad-frame} show two representative
captures from the project notes and dataset. They are useful because they
illustrate different acquisition problems.

\begin{figure}[h]
\centering
\includegraphics[width=0.9\linewidth]
  {./charuco_0000_with_auto_white_balance_and_exposure.jpg}
\caption{A ChArUco capture made with auto white balance and auto exposure.}
\label{fig:auto-frame}
\end{figure}

\begin{figure}[h]
\centering
\includegraphics[width=0.9\linewidth]{./charuco_0004.jpg}
\caption{An example of a weak calibration frame from the dataset.}
\label{fig:bad-frame}
\end{figure}

The image in Figure~\ref{fig:auto-frame} is better centered and shows more of
the board, but it still has calibration risks:

\begin{itemize}
\item the board is held by hand rather than mounted rigidly,
\item auto white balance and auto exposure may change image appearance from
      frame to frame,
\item dim indoor lighting can force longer exposure times, and
\item the board support may not remain perfectly flat and planar.
\end{itemize}

The image in Figure~\ref{fig:bad-frame} is weaker for several reasons:

\begin{itemize}
\item the board is partially cut off by the left edge of the image,
\item a large fraction of the image is empty wall rather than useful board
      coverage,
\item the board is held by hand, which increases the chance of rolling-
      shutter distortion and small motion during capture,
\item the image has a strong blue cast, suggesting that white balance is not
      well controlled, and
\item the board does not cover enough of the image to constrain the camera
      model as strongly as a larger, sharper view would.
\end{itemize}

Frames like these can still contribute some corner observations, but a strong
calibration dataset should contain many better views with the full board
visible, wider spatial coverage across the image, sharper corners, and more
stable lighting. A few weak frames are acceptable, but many weak frames can
pull the calibration toward a poor solution and inflate the reprojection
error.

\section{Inspecting Image Information from the Terminal}

It is often useful to inspect calibration images directly from the command
line before using them. This helps confirm the image format, resolution, and
basic metadata, and it can reveal whether a tool chain is preserving the
expected capture settings.

Useful commands include:

\begin{itemize}
\item \texttt{file path/to/image.jpg} to identify the file type,
\item \texttt{ffprobe -v error -show\_format -show\_streams path/to/image.jpg}
      to print image metadata in a structured way, and
\item \texttt{identify -verbose path/to/image.jpg} if ImageMagick is
      installed.
\end{itemize}

For example:

\begin{verbatim}
file ./data/charuco_left/charuco_0004.jpg
ffprobe -v error -show_format -show_streams \
  ./data/charuco_left/charuco_0004.jpg
\end{verbatim}

These terminal checks are useful when debugging a calibration dataset,
especially when image dimensions, compression, or metadata do not match the
expected camera settings.

\section{Inspecting and Changing Stream Settings from the Terminal}

When debugging webcam capture, it is also important to inspect the active
video stream itself rather than only the saved images. The main Linux tool
for this is \texttt{v4l2-ctl}.

To list the available webcams:

\begin{verbatim}
v4l2-ctl --list-devices
\end{verbatim}

To inspect one device or stream in detail:

\begin{verbatim}
v4l2-ctl -d /dev/video0 --all
v4l2-ctl -d /dev/video0 --list-ctrls-menu
\end{verbatim}

The \texttt{--all} output shows the active resolution, pixel format, frame
rate, and current control values. In our case, this kind of output revealed
that the webcam was streaming at \texttt{1280x720}, using \texttt{MJPG}, and
running with auto exposure and auto white balance enabled.

Those settings matter because they affect blur, color consistency, and the
repeatability of the calibration images.

To change camera controls, use \texttt{--set-ctrl}. For example:

\begin{verbatim}
v4l2-ctl -d /dev/video0 --set-ctrl=white_balance_automatic=0
v4l2-ctl -d /dev/video0 --set-ctrl=white_balance_temperature=4600
v4l2-ctl -d /dev/video0 --set-ctrl=auto_exposure=1
v4l2-ctl -d /dev/video0 --set-ctrl=exposure_time_absolute=157
v4l2-ctl -d /dev/video0 --set-ctrl=gain=0
v4l2-ctl -d /dev/video0 --set-ctrl=backlight_compensation=0
v4l2-ctl -d /dev/video0 --set-ctrl=power_line_frequency=2
\end{verbatim}

These commands move the webcam into a more repeatable calibration state:
manual white balance, manual exposure, low gain, disabled backlight
compensation, and 60 Hz anti-flicker mode.

Some stream parameters can also be changed directly:

\begin{verbatim}
v4l2-ctl -d /dev/video0 --set-fmt-video=width=1280,height=720,\
pixelformat=MJPG
v4l2-ctl -d /dev/video0 --set-parm=25
\end{verbatim}

After changing controls or stream settings, it is good practice to re-run
\texttt{v4l2-ctl -d /dev/video0 --all}. This confirms that the webcam
accepted the requested configuration. Not every device supports every mode,
and some controls become inactive depending on whether the camera is in
automatic or manual mode.

\section{Why We Start with Standard Calibration}

A true rolling-shutter calibration model would estimate not only intrinsics
and distortion, but also readout timing and the effect of motion during
capture. That is a more advanced problem and usually requires stronger
assumptions, more specialized solvers, or controlled motion sequences.

For \texttt{libframe}, the immediate goal is more practical:

\begin{itemize}
\item obtain stable intrinsics for each webcam,
\item calibrate the stereo geometry,
\item rectify left and right images, and
\item support downstream depth and tracking experiments.
\end{itemize}

If reprojection errors remain high after using the correct board definition
and carefully captured static images, then rolling-shutter effects become a
stronger candidate explanation. At that point, it would make sense to
investigate tighter capture discipline, synchronized acquisition, or a
dedicated rolling-shutter-aware calibration method.

\section{Summary}

Global shutter cameras match the assumptions of classical calibration
directly because each frame is captured at one instant. Rolling shutter
cameras violate that assumption when the camera or target moves during
readout. However, if motion is kept small, rolling shutter webcams can still
be calibrated effectively with a standard ChArUco-based pipeline.

The procedure in this project is therefore to use static, well-lit ChArUco
captures, estimate intrinsics and distortion with standard OpenCV
calibration, reject bad frames, and then use those results as the basis for
stereo rectification and later 3D processing.

\end{document}
